\section{Experimental Design}

This thesis proposes a survey experiment to examine how Germans form beliefs about the labor-market and social 
integration outcomes of Ukrainian and Syrian migrants. The proposed study consists of a factorial vignette experiment, 
a set of outcome measures, and a respondent questionnaire. The following sections describe each part of 
the proposed experimental design and the survey procedure.

\subsection{Target Population and Proposed Sample}

The target population consists of German citizens aged 18 or older who currently reside in Germany. Restricting 
eligibility to German citizens provides a clear operational definition of the German population whose beliefs about 
the integration of Ukrainian and Syrian migrants are the focus of this study. Similar citizenship-based sampling 
criteria have been used in previous studies examining attitudes toward refugees and migration in Germany 
\citep[e.g.,][]{hayo2023between}. Importantly, German citizenship is not used as a proxy for being native-born or 
having no migration background. Naturalized German citizens are therefore included in the target population, while 
respondents’ migration background is collected separately as a demographic characteristic in the respondent questionnaire.
At the beginning of the survey, respondents would therefore be screened with respect to age, country of residence, 
and German citizenship. Respondents who do not satisfy these eligibility criteria would not proceed to the experimental 
part of the survey.

I propose using Qualtrics to program the survey, and the online panel provider Bilendi for the recruitment of 
participants. Bilendi is suitable for the present study because it provides access to a large and diverse pool of 
respondents in Germany and allows the use of demographic quotas to approximate the composition of the target 
population. The proposed sample would use quotas for age, gender, education, and region. Respondents would receive 
compensation for their participation according to the standard compensation provided by the panel, taking into account 
the expected duration of the survey.

The final sample size would be determined based on a power analysis of the proposed experimental design. The number 
of vignette evaluations per respondent, which also affects the required sample size, is discussed in Section 3.2.

\subsection{Factorial Vignette Design}

The study employs a factorial vignette experiment in which respondents evaluate hypothetical migrants whose 
characteristics are independently varied across profiles. This approach is particularly suitable for the present 
research question because migrants differ simultaneously along several characteristics. By independently varying these 
characteristics, the design makes it possible to distinguish their effects on respondents’ beliefs while holding the remaining 
characteristics constant. The proposed design builds on previous conjoint and factorial vignette experiments examining 
how migrant characteristics shape evaluations \citep{hainmueller2015hidden, bansak2016economic, czymara2016deutschland}.

The proposed vignettes describe migrants who have come to Germany and are currently looking for stable employment. 
Each profile varies independently along five binary dimensions: gender, country of origin, reason for migration, 
human capital, and German-language proficiency. 
These vignette attributes are selected to capture migrant characteristics that are central to the research questions 
and that previous research has shown to influence perceptions of migrants. Country of origin is varied between 
(1) Ukrainian and (2) Syrian migrants and represents the central comparison of the study. Reason for migration 
distinguishes between the two types of migrants: (1) refugees who fled the wars in Ukraine or Syria and (2) economic 
migrants who came to Germany in search of better employment opportunities. Educational attainment captures 
differences in migrants' human capital, which previous vignette and conjoint studies have shown to influence 
evaluations of immigrants \citep{hainmueller2010attitudes, hainmueller2015hidden, czymara2016deutschland}. 
The design distinguishes between (1) migrants without a formal vocational qualification and (2) migrants with a 
university degree. German-language proficiency is varied between (1) limited and (2) good German skills. 
Language proficiency may be relevant both for social integration and for labor-market outcomes, since knowledge of the 
host-country language represents an important form of human capital \citep{czymara2016deutschland}. 
Finally, gender distinguishes between (1) male and (2) female migrants. Gender is included because previous 
experimental research suggests that migrant gender can affect how immigrants and refugees are perceived, including in 
the German context \citep{czymara2017refugees, gereke2020gendered}. Randomizing gender therefore allows the effects 
of the other vignette characteristics to be estimated across both male and female migrant profiles.

As discussed previously, existing research shows that religion and perceived cultural distance can influence attitudes 
toward immigrants \citep[e.g.,][]{czymara2016deutschland}. However, in this study, religious affiliation is 
deliberately left unspecified. The main interest of this thesis is in the overall beliefs that respondents associate 
with Ukrainian and Syrian migrants. Explicitly including religion as an additional vignette attribute would not only 
make the design more complex, but could also change the beliefs that respondents would otherwise associate with the 
migrant's country of origin. For this reason, religion is not varied separately in the experiment. The estimated 
country-of-origin effect should therefore be interpreted as an overall effect of Ukrainian versus Syrian origin, 
which may also reflect perceived cultural or religious differences between the two groups.

Table \ref{tab:vignette_attributes} summarizes the five vignette attributes and their respective levels and shows how 
the randomized characteristics relate to the hypotheses developed in Chapter 2.
\begin{table}[htbp]
\centering
\caption{Vignette Attributes and Attribute Levels}
\label{tab:vignette_attributes}

\renewcommand{\arraystretch}{1.25}

\begin{tabular}{p{3.1cm} p{3.7cm} p{3.0cm} l}
\toprule
\textbf{\makecell[l]{Vignette\\attribute}} &
\textbf{Level 1} &
\textbf{Level 2} &
\textbf{Hypothesis} \\
\midrule

Country of origin
    & Syria
    & Ukraine
    & H3a, H3b, H3c \\

\addlinespace
\makecell[l]{Reason for\\migration}
    & Economic migrant
    & War refugee
    & H2, H3c \\

\addlinespace
Human capital
    & \makecell[l]{No formal vocational\\qualification}
    & \makecell[l]{University\\degree}
    & H1a \\

\addlinespace
\makecell[l]{German-language\\proficiency}
    & Limited
    & Good
    & H1b \\

\addlinespace
Gender
    & Male
    & Female
    & -- \\

\bottomrule
\end{tabular}

\vspace{0.3em}

\begin{minipage}{0.95\textwidth}
\footnotesize
\textit{Notes:} The table presents the five independently randomized
attributes used to construct the migrant profiles. Gender is included
as an additional randomized characteristic but is not associated with
a separate hypothesis.
\end{minipage}

\end{table}

The five binary attributes generate a full factorial design of \(2^5=32\) possible migrant profiles.
The vignettes would be worded as follows:\bigskip

[Mr./Mrs.] [random letter]. came to Germany from [Ukraine/Syria]. [He/She] came to Germany as [a refugee after fleeing 
the war in his/her country / an economic migrant in search of better employment and economic opportunities in Germany]. 
[He/She] [has no formal vocational qualification / has completed a university degree], speaks German 
[only to a limited extent / well], and is currently looking for stable employment.\bigskip

Since evaluating 
all possible profiles would place a substantial burden on respondents, I propose that each respondent evaluate a 
randomly selected subset of six vignettes. This provides several profile evaluations per respondent while keeping 
the survey manageable, particularly because each vignette is followed by several belief questions about migrant outcomes. 
The choice therefore reflects a trade-off between obtaining sufficient statistical information and limiting 
cognitive burden and respondent fatigue, which are important considerations in factorial survey designs 
\citep{sauer2011application}. The exact number of vignettes could be further assessed through pilot testing prior 
to implementation.


\subsection{Outcome Variables}
The five outcome questions are going to be answered right after each vignette. They are designed to capture the respondents'
beliefs about some specific labor-market and social integration outcomes of the migrant based on the profile they just saw by
evaluating it according to the vignette attributes. 

The outcomes cover three aspects: beliefs about labor-market integration outcomes (expected employment and expected income), 
beliefs about labor-market discrimination (expected hiring and workplace treatment), and beliefs about social 
integration. The exact wording of all outcome questions and response scales is provided in Appendix [X]. 

The first two outcome questions capture beliefs about the migrant's labor-market integration. 
As discussed in Section 2.2, labor-market integration is reflected in several dimensions, including employment, 
earnings, and occupational outcomes \citep{brell2020labor, fasani2024refugees}. Therefore, to capture the clearest
and the most interpretable dimensions, the proposed experiment measures respondents' beliefs about labor-market 
integration based on expected employment and expected income.
First, respondents state the probability that the migrant will find regular paid employment subject to social security 
contributions within the next two years. The type of employment in this question is formulated in a way to exclude 
informal, marginal, or unstable jobs. The timeframe of two years is chosen because it is long enough for job search, 
language learning, and recognition of qualifications of the migrants in Germany, but still short enough for respondents 
to imagine. This question will be answered on a scale from 0 to 100 percent. This probability scale provides a direct 
quantitative measure of respondents' beliefs about employment prospects. 
Third, respondents indicate the migrant's expected monthly net income five years from now by selecting one of seven 
income categories.

Two additional measures capture respondents' beliefs about discrimination in the German labor market. The first concerns 
hiring discrimination and asks respondents how likely is it that the migrant would be treated less favorably in the 
hiring process than an otherwise comparable German applicant. The second concerns discrimination after entering 
employment and measures the expected likelihood of unequal treatment in the workplace, for example in areas such as 
wages, promotion opportunities, or interactions with colleagues. Both outcomes are measured on seven-point Likert scales, 
with higher values indicating greater agreement and therefore higher expected discrimination.

The final outcome captures beliefs about social integration. Respondents assess how likely the migrant is to 
successfully integrate socially in Germany within the next five years. To give respondents a common understanding of 
social integration, the question defines it as the migrants building regular social contacts with German citizens within
the next five years. Responses are recorded on a seven-point Likert scale, with higher values indicating more favorable 
beliefs about social integration.

Combined, these five measures capture different dimensions of respondents' beliefs about migrant integration rather 
than a single overall evaluation. The labor-market variables provide the main outcomes for H1a, H1b, H2, and H3a, 
while the discrimination measures provide the outcomes for H3b and the social-integration measure provides an additional 
outcome for H1b and H3a. All five outcome measures are used to examine H3c. In the empirical 
analysis, the outcome measures are analyzed separately rather than being combined into a single integration index.

\subsection{Respondent Questionnaire}

\subsection{Randomization and Survey Procedure}